\chapter[Discusión y conclusiones]{Discusión y\\ conclusiones}
\label{chap:conclusiones}

\section{Interpretación de los resultados}

La comparación desarrollada en este trabajo permite valorar cómo se aprovechan unos pocos ejemplos etiquetados al ajustar una CNN o al incorporarlos al contexto de Gemma~4. El contraste entre los resultados del simulador y los del HUCA muestra que su utilidad depende también de cuánto se corresponden con los casos que se pretende clasificar.

\subsection{Aprendizaje con datos del simulador}

Ambos métodos mejoran al aumentar \(K\) en el simulador y alcanzan sus mejores resultados con 32 ejemplos. La CNN mejora de forma continua mientras que Gemma~4 llega a un rendimiento comparable y obtiene una ligera ventaja con la selección balanceada en ese último punto. Esta evolución responde a la primera pregunta del trabajo y es coherente con un aprendizaje favorecido por la homogeneidad de las imágenes y por una relación directa entre la forma del trazado y la etiqueta.

Los latidos comparten el mismo formato y varían dentro de las formas previstas por el generador. Las regiones del ECG resultan por ello más fáciles de comparar entre ejemplos y las etiquetas corresponden a rasgos presentes en la propia imagen. Aumentar los ejemplos permite mostrar más variaciones de esas combinaciones sin introducir toda la diversidad de los registros clínicos. El buen rendimiento indica que ambos modelos pueden aprovechar esa regularidad aunque no permite asegurar que hayan aprendido únicamente los rasgos previstos ni que reconozcan todas las manifestaciones clínicas de Brugada.

La dependencia de Gemma~4 respecto de las demostraciones refuerza esta interpretación porque no detecta positivos cuando recibe solo las instrucciones. Los ejemplos parecen ayudarle a concretar qué diferencias visuales debe asociar con cada respuesta y esa orientación resulta suficiente para clasificar otros latidos del mismo generador. Su éxito en este escenario debe entenderse como aprendizaje dentro de una tarea controlada cuya utilidad clínica depende únicamente de lo que se conserve al pasar a pacientes reales.

\subsection{Adaptación al dominio de los datos clínicos}

La evaluación sobre el HUCA muestra que se conserva muy poco de ese rendimiento y que añadir ejemplos del simulador no resuelve la transferencia. La CNN mejora al reconocer los latidos generados mientras pierde rendimiento con los pacientes reales y Gemma~4 con demostraciones sintéticas sigue clasificando casi todos los ECG clínicos como negativos. Este contraste es compatible con un aprendizaje de regularidades propias del generador que dejan de ser suficientes cuando cambian las imágenes y la tarea que representan sus etiquetas.

La homogeneidad que facilita el aprendizaje en el simulador limita también la variedad de formas que los modelos conocen a través de sus ejemplos. Los registros clínicos presentan más variación y su etiqueta corresponde al paciente en lugar de describir directamente la combinación de rasgos usada en el simulador. La pérdida de rendimiento puede deberse tanto a la apariencia del trazado como a esa diferencia entre las tareas.

Las variantes de adaptación intentan reducir las diferencias visuales pero siguen aprendiendo con las etiquetas del simulador. CORAL y MMD mejoran frente a la CNN base con el mismo \(K\) aunque ninguna de las variantes supera el mejor resultado sin adaptación y el rendimiento permanece próximo al azar. El experimento no separa el efecto de una alineación visual insuficiente del cambio de tarea pero sí muestra que la adaptación probada no compensa ambas dificultades y que los ejemplos sintéticos no son capaces de aportar el contexto suficiente para esta clasificación.

\subsection{Aprendizaje con datos clínicos}

El origen de los ejemplos cambia el resultado cuando Gemma~4 recibe demostraciones del HUCA y clasifica otros pacientes del mismo conjunto. Su rendimiento mejora al aumentar \(K\) y supera al de la CNN entrenada con esos datos desde ocho ejemplos. Con dos y cuatro ambos métodos permanecen cerca del azar aunque la CNN obtiene una ligera ventaja con apoyo de validación etiquetada adicional. La comparación responde así a la segunda pregunta del trabajo y muestra que el aprendizaje en contexto aprovecha mejor los ejemplos clínicos en los tamaños mayores evaluados.

Las demostraciones del HUCA comparten con la consulta tanto la variación real del trazado como la definición de la etiqueta que se quiere predecir. Esta correspondencia ayuda a explicar la mejora respecto del uso de ejemplos sintéticos y sugiere que Gemma~4 puede utilizar sus características visuales previas para relacionar nuevos casos con las referencias del contexto. Para la CNN evaluada resulta imposible modelar adecuadamente la patología con tan pocos ejemplos.

La mejora en el dominio clínico sigue siendo parcial porque incluso la mejor configuración omite aproximadamente una cuarta parte de los positivos. Aunque no alcanza valores comparables a los del simulador Gemma~4 podría servir como herramienta de triaje para priorizar la revisión de casos sospechosos y facilitar un estudio más exhaustivo de la enfermedad.

\section{Aportaciones y limitaciones}
\label{sec:limitaciones}

La principal aportación del trabajo es una comparación de dos formas de aprender con pocos ejemplos que contempla tanto un entorno controlado como su relación con datos clínicos. La evaluación por paciente permite estudiar el efecto de \(K\) sin incluir al paciente de prueba entre los ejemplos etiquetados de cada pliegue y la separación de escenarios ayuda a distinguir el reconocimiento de las familias sintéticas de la clasificación con etiquetas del HUCA. El simulador aporta un marco para estudiar combinaciones morfológicas conocidas y la comparación con pacientes reales permite identificar hasta dónde llegan esas reglas al enfrentarse a la tarea clínica.

El alcance de los resultados está limitado por la falta de anotaciones morfológicas por derivación con las que comprobar si los modelos reconocen los rasgos que motivan la sospecha de Brugada o si aprovechan otras regularidades de las imágenes. Esta limitación afecta especialmente a la interpretación de una etiqueta asignada al paciente a partir de información que puede exceder lo visible en el latido V1 seleccionado.

La muestra procede de un único centro hospitalario y excluye los casos cuya clasificación no estaba clara al preparar el estudio. Su composición no representa necesariamente la variedad de pacientes que llegaría a un circuito de triaje y limita la posibilidad de trasladar directamente las métricas a ese contexto. El uso de una sola derivación y de un latido por paciente también reduce la información disponible frente a una valoración clínica completa que puede requerir otras derivaciones o registros obtenidos en distintos momentos.

La falta de validación externa y prospectiva impide valorar todavía la utilidad en el triaje de las configuraciones estudiadas. No se han medido la calibración de las probabilidades ni el efecto de los errores sobre la carga de revisión y tampoco se dispone de una evaluación suficientemente precisa por subgrupos.

\section{Continuación del trabajo}

La continuación del proyecto debería comenzar con la revisión especializada de las etiquetas y de los casos dudosos para disponer de una base clínica contrastada sobre la que repetir la evaluación. Incorporar anotaciones de los rasgos visibles en cada derivación ayudaría a relacionar la respuesta de los modelos con la morfología del ECG y a distinguir los errores de reconocimiento de las discordancias entre una imagen aislada y la etiqueta del paciente. Esa información también permitiría revisar el simulador para acercar sus reglas a una tarea clínica definida con mayor precisión.

La ampliación del número de pacientes y de casos positivos permitiría estudiar si la CNN mejora de forma sostenida al disponer de más datos. Sobre esa base tendría sentido comparar el ajuste completo de ResNet18 con un extractor de características congelado y con arquitecturas especializadas en ECG. La incorporación de otras derivaciones y de varios latidos por paciente ayudaría además a comprobar cuánto limita el uso exclusivo de V1 y si una representación más completa aporta información útil para la clasificación.

También podría explorarse el ajuste fino de Gemma~4 con datos del simulador como etapa de adaptación al análisis de imágenes de ECG. Generar más ejemplos con mayor variedad de trazados permitiría actualizar sus pesos para reforzar el reconocimiento de los rasgos morfológicos y combinar después ese ajuste con demostraciones clínicas. Su evaluación sobre el HUCA permitiría comprobar si esa preparación mejora la clasificación respecto del modelo original o si el aprendizaje continúa ligado a las regularidades del simulador.

La continuación del estudio del aprendizaje en contexto pasa por estudiar la selección de demostraciones por similitud y el efecto de su orden dentro de la consulta. Mantener controlados esos factores permitiría valorar mejor la contribución del balance de clases y comprobar si la ventaja observada con Gemma~4 se conserva con otros modelos o instrucciones.

La evaluación con pacientes de otros hospitales sería necesaria para comprobar la generalización frente a cambios de equipo o de presentación de las imágenes. Ese estudio debería reservar datos ajenos a cualquier ajuste y acompañarse de intervalos de incertidumbre y de una revisión experta de los errores. Una fase prospectiva posterior permitiría medir el efecto sobre la priorización de casos y sobre la carga de trabajo antes de valorar una aplicación de triaje clínico.

La posible utilización de Gemma~4 para apoyar el triaje debe estudiarse dentro de ese proceso de validación. Priorizar ECG para su revisión podría facilitar la recogida de nuevos casos confirmados con los que investigar después modelos supervisados de mayor rendimiento o con mayor explicabilidad.

\section{Conclusiones}

\begingroup
\widowpenalty=10000
Este Trabajo de Fin de Grado ha comparado el ajuste supervisado de una CNN ResNet18 y el aprendizaje en contexto de Gemma~4 para clasificar imágenes de ECG con pocos ejemplos etiquetados. La evaluación muestra que ambos métodos pueden distinguir las clases del simulador y mejorar al aumentar los ejemplos aunque ese aprendizaje no se traslada de forma suficiente a los pacientes del HUCA. Las mejoras limitadas de las variantes de adaptación confirman que las soluciones probadas no resuelven la diferencia entre el entorno generado y la tarea definida por las etiquetas clínicas.
\par
\endgroup

La aportación más relevante para el escenario con pocos datos es la mejora de Gemma~4 cuando recibe demostraciones del propio HUCA. Su ventaja frente a la CNN entrenada con esos datos desde ocho ejemplos muestra que el aprendizaje en contexto puede resultar útil para clasificar otros pacientes del mismo conjunto sin ajustar los pesos del modelo. La ausencia de detección de positivos sin demostraciones y la escasa transferencia desde el simulador indican que esa utilidad depende de ofrecer ejemplos pertinentes para la tarea y no puede deducirse únicamente del conocimiento previo del modelo.

El trabajo permite así identificar una vía de investigación para el apoyo al triaje y precisar las condiciones que todavía deben estudiarse. Los resultados corresponden a la muestra y a las configuraciones evaluadas y necesitan contrastarse mediante revisión clínica de las etiquetas y validación en otros centros. Su posible aplicación se orienta a priorizar la revisión especializada de casos sospechosos dentro de la evaluación clínica completa. Comprobar que esa priorización mejora la atención sin excluir pacientes que requieren valoración es el siguiente paso para determinar su utilidad.
